\documentclass[11pt,a4paper]{ctexart}
\usepackage[a4paper,margin=1.60cm,top=1.8cm,bottom=1.8cm]{geometry}
\usepackage{ipara-reading}

% ==============================================================================
% 单篇精读训练手札模板
%
% 文件职责：
% 1. article.md       = 原文阅读副本、来源信息与学习者首读标记；
% 2. reading_view.tex = 本篇文章的精读、声音与输出训练现场；
% 3. 02_expressions/  = 跨文章表达节点的长期维护位置；
% 4. 03_capabilities/ = 跨文章能力状态的长期维护位置。
%
% TeX 可以为了 Cornell 排版再次呈现原文，但不要复制表达层的完整节点，
% 也不要在单篇手札中维护跨文章能力结论。各区块按文章价值和真实训练需要增删。
% ==============================================================================

\begin{document}

% ------------------------------------------------------------------------------
% 1. 题头：只保留识别文章和训练所需的信息
% ------------------------------------------------------------------------------
\articleheader
  {【文章主标题】}
  {【可选副标题】}
  {【杂志 / 网站 / 栏目】}
  {【主题】}
  {【可选：训练标签】}
  {【原始发布日期或训练日期】}

% ------------------------------------------------------------------------------
% 2. 双栏精读：优先从 article.md 的学习者首读标记出发；右栏只解决当前真实问题
% ------------------------------------------------------------------------------
\readingsection{一、双栏精读}

\begin{paracol}{2}

% ---------------- 第 1 段 ----------------
\artpara{P1}{
  【粘贴第 1 段原文。只标真正需要处理的位置，例如
  \obs{真实理解断点}{1}；若某个表达值得继续观察，可用
  \kw{候选表达}{2} 轻量标出。】
}

\switchcolumn

\noteobs{1}{真正的理解断点}
  {先判断卡点属于词义、搭配、句法挂接、指代还是背景知识，再做定点解释。不要用整段翻译替代诊断。}

\noteitem{2}{候选表达 · 若已入库则标注 EXP-xxx}
  {只写理解当前语境所需的简短含义。}
  {只有在口语重构确实有帮助时，才补一条自然说法。}
  {若表达已进入长期表达层，在这里给出编号或去向；否则保留为本篇候选即可。}

\switchcolumn*

% ---------------- 第 2 段 ----------------
\artpara{P2}{
  【粘贴下一段原文。只有真实存在理解、表达或声音价值时才增加右栏笔记。】
}

\switchcolumn

% 没有值得记录的断点时，右栏可以保持很短，甚至不添加条目。

\end{paracol}

% ------------------------------------------------------------------------------
% 3. 训练后再还原文章结构，不在 article.md 中提前给答案
% ------------------------------------------------------------------------------
\readingsection{二、文章结构与主旨复原}

\begin{coretakeaway}[训练后主旨]
  训练后，用自己的英语概括文章最核心的意思。不要照抄导语；如果仍说不清，先回到文章结构，不急着做语言润色。
\end{coretakeaway}

\begin{itemize}
  \item \textbf{结构节点 1}：根据文章真实文体填写，例如场景、问题或核心发现。
  \item \textbf{结构节点 2}：记录文章如何推进，例如证据、原因、经历或转折。
  \item \textbf{结构节点 3}：记录文章如何收束，例如不同观点、限制、结果或判断。
\end{itemize}

% ------------------------------------------------------------------------------
% 4. 语音与韵律：先有可靠声音参考，再用文本标注意群、重音与语调；数量自适应
% ------------------------------------------------------------------------------
\readingsection{三、语音与韵律训练}

\begin{prosodybox}[黄金句 · 示例]
  \prosodysource{【填写真人原音、词典音频或 TTS 中性参考；有时间戳时一并记录。若只有文字，明确写“仅文本参考”。】}
  \prosodyorig{【粘贴一条真正值得反复模仿的原句。】}
  \prosodyipa{\textbf{关键词 1} \ipa{/ˈ.../} \quad$\bm{\cdot}$\quad \textbf{关键词 2} \ipa{/ˈ.../}}
  \prosodychunks{【按真实表达意图标出 \stress{信息重音} \chunk 与自然意群 \chunk；必要时补充语调走向。】 \tonefall}
  \prosodyfocus{【列出本句真正承担新信息、对比或核心动作的焦点词。】}
  \prosodyflaws{
    \item 【只记录真实练习后仍需要提醒的发音、弱读、衔接或节奏问题。】
  }
\end{prosodybox}

% 根据文章价值增删 prosodybox；不为凑数设置固定配额。

% ------------------------------------------------------------------------------
% 5. 输出训练：先保存真实首答，再给参考复述与差异对照
% ------------------------------------------------------------------------------
\readingsection{四、脱稿输出与对照}

\oralblock{第一次真实概括 / 复述}{
  【不看原文，用自己的英语讲清文章核心内容。原样保留学习者版本；若来自语音转写，只按内容与句型分析，不据此判断发音、停顿或真实流利度。】
}

\oralblock{关键诊断}{
  【只记录最影响收益的少量差异：内容准确性、结构组织、语言可调用性。个人观点与原文事实要明确区分。】
}

\oralblock{AI 参考复述}{
  【在学习者已经完成真实输出并接受关键反馈后再填写。事实只来自原文，结构清楚，语言自然但不过度拔高，难度约为学习者当前输出 +1。必须明确这是 AI 参考版，不能当作学习者已经会说的证据。】
}

\oralblock{对照：我的版本 vs 参考版}{
  【简短记录四类差异：已经说出来的内容、需要修正的事实、值得带走的 2--4 个语块、参考版更稳定的结构。不要做逐句红笔批改。】
}

\oralblock{个人观点 / 迁移讨论}{
  【选择真正值得讨论的问题。重点观察观点能否展开、表达能否自然调用，不追求表面上的“高级词汇”。】
}

\oralblock{下一次冷检索}{
  【写清下一次先不看参考答案要复原什么，以及无提示检查哪 1--2 个表达。】
}

% ------------------------------------------------------------------------------
% 6. AI 对谈：只保留关键片段；跨文章稳定问题由能力层维护
% ------------------------------------------------------------------------------
\readingsection{五、对谈与关键反馈}

\begin{aiinteraction}{【本篇最值得继续谈的问题】}
  \aiquestion{【填写一个能够引出真实观点、而不是只复述原文的问题。】}
  \aiuseranswer{【保留学习者当时的真实说法；没有真实输出时留空，不由 AI 代填。】}
  \aicorrection{
    \begin{itemize}
      \item 只保留最影响理解、自然度或目标表达调用的关键反馈。
    \end{itemize}
  }
  \ainatural{【只有确实能改善表达时才给自然改写；若其中产生新的长期表达，再链接到表达层。】}
\end{aiinteraction}

% ------------------------------------------------------------------------------
% 7. 本篇向外沉淀什么：这里只列引用，不复制长期资产
% ------------------------------------------------------------------------------
\readingsection{六、本篇沉淀与跨层引用}

\begin{itemize}
  \item \textbf{\color{readGreen}进入表达层}：\texttt{EXP-xxx} · 【表达名称】；\texttt{EXP-yyy} · 【表达名称】。
  \item \textbf{\color{readMain}仍留本篇}：只在当前主题中有用的术语、一次性背景或局部修辞。
  \item \textbf{\color{readGray}能力证据}：若本次训练构成代表性证据，记录能力文件与证据锚点；否则只保留为本篇现象。
  \item \textbf{\color{readGray}后续复习}：记录下一次先不看 AI 参考复述要冷检索的主旨、结构和 1--2 个表达，不按固定数量凑项。
\end{itemize}

\end{document}
